\scanpage{24}
\mspage{3}
liest Deutsch. Vielleicht könnte ich all das, was irgendwie direkt ``autobiographisch'' ist, auf Französisch schreiben -- doch eingehende Meditationen, \unsure{insbesondére|insbesondere} an Hand von deutschsprachigem Material, auf Französisch?

Die Möglichkeit fiel mir auch ein, die beiden obigen Ansätze simultan durchzuführen -- nämlich eine \unsure{Chronologische|chronologische} Untersuchung über mein Leben \inserted{sofort} anzufangen, und zu gleicher Zeit die Lektüre etwa meines Briefwechsels mit \unsure{Hanka}.

zu befürchten wär, dass ich dauernd dabei wäre, nicht vorhandene Erinnerungen zusammenzukratzen zu versuchen! Von dort aus würden \deleted{sich} sicher viele Kernfragen \deleted{\unsure{behandelt|beantworten}} \inserted{auftauchen} und \deleted{\unsure{be}} \deleted{\unsure{behandelt}} \inserted{untersucht} werden, möglicherweise auch andere Perioden \deleted{\unsure{sichtbar}} \inserted{zur} Untersuchung aufdrängen. Wenn die Untersuchung der \unsure{Kernperio|Kernperiode} oder des Ausgangsmaterials fürs erste beendet wäre, wäre die Zeit vielleicht ``reifer'', um mich an eine chronologische Folge zu halten. Jedenfalls
% archive sheet number: 23
